% ========================================================
% BÀI TẬP TOÁN 5 - KẾT NỐI TRI THỨC VỚI CUỘC SỐNG
% Bài 4: Phân số thập phân
% Định dạng: BTTeX Part III - Trắc nghiệm điền số ngắn
% ========================================================

% Câu 1
\begin{ex}
  Viết phân số $\frac{3}{5}$ thành phân số thập phân có mẫu số là $10$. Tử số của phân số thập phân đó là bao nhiêu?
  \shortans{6}
  \loigiai{
    Nhân cả tử và mẫu với 2: $\frac{3 \times 2}{5 \times 2} = \frac{6}{10}$. Tử số là 6.
  }
\end{ex}

% Câu 2
\begin{ex}
  Viết phân số $\frac{7}{20}$ thành phân số thập phân có mẫu số là $100$. Tử số của phân số đó là bao nhiêu?
  \shortans{35}
  \loigiai{
    Nhân cả tử và mẫu với 5: $\frac{7 \times 5}{20 \times 5} = \frac{35}{100}$. Tử số là 35.
  }
\end{ex}

% Câu 3
\begin{ex}
  Viết phân số $\frac{9}{25}$ thành phân số thập phân có mẫu số là $100$. Tử số của phân số đó là bao nhiêu?
  \shortans{36}
  \loigiai{
    Nhân cả tử và mẫu với 4: $\frac{9 \times 4}{25 \times 4} = \frac{36}{100}$. Tử số là 36.
  }
\end{ex}

% Câu 4
\begin{ex}
  Viết phân số $\frac{3}{125}$ thành phân số thập phân có mẫu số là $1\,000$. Tử số của phân số đó là bao nhiêu?
  \shortans{24}
  \loigiai{
    Nhân cả tử và mẫu với 8: $\frac{3 \times 8}{125 \times 8} = \frac{24}{1000}$. Tử số là 24.
  }
\end{ex}

% Câu 5
\begin{ex}
  Rút gọn phân số $\frac{45}{300}$ thành phân số thập phân có mẫu số là $100$. Tử số của phân số đó là bao nhiêu?
  \shortans{15}
  \loigiai{
    Chia cả tử và mẫu cho 3: $\frac{45 : 3}{300 : 3} = \frac{15}{100}$. Tử số là 15.
  }
\end{ex}

% Câu 6
\begin{ex}
  Rút gọn phân số $\frac{64}{800}$ thành phân số thập phân có mẫu số là $100$. Tử số của phân số đó là bao nhiêu?
  \shortans{8}
  \loigiai{
    Chia cả tử và mẫu cho 8: $\frac{64 : 8}{800 : 8} = \frac{8}{100}$. Tử số là 8.
  }
\end{ex}

% Câu 7
\begin{ex}
  Chuyển phân số $\frac{13}{2}$ thành phân số thập phân có mẫu số là $10$. Tử số của phân số đó là bao nhiêu?
  \shortans{65}
  \loigiai{
    Nhân cả tử và mẫu với 5: $\frac{13 \times 5}{2 \times 5} = \frac{65}{10}$. Tử số là 65.
  }
\end{ex}

% Câu 8
\begin{ex}
  Viết phân số $\frac{21}{250}$ thành phân số thập phân có mẫu số là $1\,000$. Tử số của phân số đó là bao nhiêu?
  \shortans{84}
  \loigiai{
    Nhân cả tử và mẫu với 4: $\frac{21 \times 4}{250 \times 4} = \frac{84}{1000}$. Tử số là 84.
  }
\end{ex}

% Câu 9
\begin{ex}
  Tìm số tự nhiên $x$ để $\frac{x}{100} = \frac{3}{4}$.
  \shortans{75}
  \loigiai{
    Ta có $\frac{3}{4} = \frac{3 \times 25}{4 \times 25} = \frac{75}{100}$. Do đó $x = 75$.
  }
\end{ex}

% Câu 10
\begin{ex}
  Tìm số tự nhiên $y$ để $\frac{12}{50} = \frac{y}{100}$.
  \shortans{24}
  \loigiai{
    $\frac{12}{50} = \frac{12 \times 2}{50 \times 2} = \frac{24}{100}$. Do đó $y = 24$.
  }
\end{ex}

% Câu 11
\begin{ex}
  Trong các phân số $\frac{3}{10}, \frac{15}{100}, \frac{7}{25}, \frac{9}{1000}$, có bao nhiêu phân số là phân số thập phân?
  \shortans{3}
  \loigiai{
    Các phân số thập phân là các phân số có mẫu số là 10, 100, 1000,... Ở đây có 3 phân số thỏa mãn là $\frac{3}{10}, \frac{15}{100}, \frac{9}{1000}$.
  }
\end{ex}

% Câu 12
\begin{ex}
  Tính tổng của hai phân số thập phân: $\frac{3}{10} + \frac{45}{100}$. Kết quả viết dưới dạng phân số thập phân có mẫu số $100$. Tử số là bao nhiêu?
  \shortans{75}
  \loigiai{
    Quy đồng mẫu 100: $\frac{30}{100} + \frac{45}{100} = \frac{75}{100}$. Tử số là 75.
  }
\end{ex}

% Câu 13
\begin{ex}
  Tính hiệu của hai phân số thập phân: $\frac{8}{10} - \frac{25}{100}$. Kết quả viết dưới dạng phân số thập phân có mẫu số $100$. Tử số là bao nhiêu?
  \shortans{55}
  \loigiai{
    $\frac{80}{100} - \frac{25}{100} = \frac{55}{100}$. Tử số là 55.
  }
\end{ex}

% Câu 14
\begin{ex}
  Tính giá trị biểu thức: $\frac{7}{10} + \frac{23}{100} + \frac{70}{1000}$. Kết quả viết dưới dạng phân số thập phân mẫu số $1\,000$. Tử số là bao nhiêu?
  \shortans{1000}
  \loigiai{
    $\frac{700}{1000} + \frac{230}{1000} + \frac{70}{1000} = \frac{1000}{1000}$. Tử số là 1000.
  }
\end{ex}

% Câu 15
\begin{ex}
  Một mảnh vườn có diện tích là $100\text{ m}^2$, trong đó người ta dùng $\frac{35}{100}$ diện tích để trồng hoa. Diện tích trồng hoa là bao nhiêu mét vuông?
  \shortans{35}
  \loigiai{
    Diện tích trồng hoa là: $100 \times \frac{35}{100} = 35\text{ m}^2$.
  }
\end{ex}

% Câu 16
\begin{ex}
  Một kho lương thực có $1\,000$ tấn thóc, đã xuất bán $\frac{275}{1000}$ số thóc. Hỏi kho đó đã xuất bán bao nhiêu tấn thóc?
  \shortans{275}
  \loigiai{
    Số thóc xuất bán là: $1\,000 \times \frac{275}{1000} = 275$ tấn.
  }
\end{ex}

% Câu 17
\begin{ex}
  Tìm phân số thập phân có mẫu số là $10$ nằm giữa $\frac{1}{2}$ và $\frac{4}{5}$. Tử số của phân số đó là bao nhiêu?
  \shortans{7}
  \loigiai{
    Quy đồng mẫu 10: $\frac{5}{10} < \frac{x}{10} < \frac{8}{10}$. Phân số $\frac{7}{10}$ nằm giữa $\frac{5}{10}$ và $\frac{8}{10}$. Tử số là 7.
  }
\end{ex}

% Câu 18
\begin{ex}
  Viết số $4$ dưới dạng phân số thập phân có mẫu số là $100$. Tử số là bao nhiêu?
  \shortans{400}
  \loigiai{
    $4 = \frac{4 \times 100}{100} = \frac{400}{100}$. Tử số là 400.
  }
\end{ex}

% Câu 19
\begin{ex}
  Phân số $\frac{75}{100}$ sau khi rút gọn về phân số tối giản có tổng của tử số và mẫu số bằng bao nhiêu?
  \shortans{7}
  \loigiai{
    $\frac{75}{100} = \frac{3}{4}$. Tổng tử số và mẫu số là $3 + 4 = 7$.
  }
\end{ex}

% Câu 20
\begin{ex}
  Phân số $\frac{125}{1000}$ khi rút gọn về phân số tối giản có mẫu số bằng bao nhiêu?
  \shortans{8}
  \loigiai{
    Chia cả tử và mẫu cho 125: $\frac{125 : 125}{1000 : 125} = \frac{1}{8}$. Mẫu số là 8.
  }
\end{ex}

% Câu 21
\begin{ex}
  Tìm số tự nhiên $a$ biết: $\frac{a}{10} + \frac{a}{10} = \frac{6}{10}$.
  \shortans{3}
  \loigiai{
    $\frac{2a}{10} = \frac{6}{10} \implies 2a = 6 \implies a = 3$.
  }
\end{ex}

% Câu 22
\begin{ex}
  Một cửa hàng có $100$ hộp bánh, buổi sáng bán được $\frac{3}{10}$ số hộp bánh, buổi chiều bán được $\frac{45}{100}$ số hộp bánh. Cả ngày cửa hàng bán được bao nhiêu hộp bánh?
  \shortans{75}
  \loigiai{
    Buổi sáng bán được $100 \times \frac{3}{10} = 30$ hộp. Buổi chiều bán được $100 \times \frac{45}{100} = 45$ hộp. Cả ngày bán được $30 + 45 = 75$ hộp.
  }
\end{ex}

% Câu 23
\begin{ex}
  Viết phân số $\frac{1}{8}$ thành phân số thập phân có mẫu số là $1\,000$. Tử số của phân số đó là bao nhiêu?
  \shortans{125}
  \loigiai{
    $\frac{1 \times 125}{8 \times 125} = \frac{125}{1000}$. Tử số là 125.
  }
\end{ex}

% Câu 24
\begin{ex}
  Tìm chữ số $x$ để $\frac{4x}{100} = \frac{2}{5}$.
  \shortans{0}
  \loigiai{
    Ta có $\frac{2}{5} = \frac{40}{100}$. Suy ra số có hai chữ số là 40, do đó $x = 0$.
  }
\end{ex}

% Câu 25
\begin{ex}
  Có bao nhiêu phân số thập phân có mẫu số bằng $100$ và lớn hơn $\frac{1}{10}$ nhưng bé hơn $\frac{2}{10}$?
  \shortans{9}
  \loigiai{
    Quy đồng mẫu 100: $\frac{10}{100} < \frac{x}{100} < \frac{20}{100}$. Số giá trị của $x$ là từ 11 đến 19: có $19 - 11 + 1 = 9$ số.
  }
\end{ex}
